\section{Experimental setting}
\label{sec:setting}

\paragraph{Data.}
We use participant S1 (\texttt{sub-01}) of Chisco \citep{zhang2024chisco}, distributed as OpenNeuro
\texttt{ds005170} \citep{ds005170}. In each trial the participant reads and memorises a sentence (the first 5\,s), then imagines
saying it (the next 3.3\,s). The paper reports 6{,}681 presented sentences per
participant over five recording days. The released preprocessed imagined-speech epochs for S1
contain 6{,}243 of them: within each of the 45 runs, the released trials are an in-order subsequence
of that run's stimulus list, so preprocessing removed 438 trials and reordered none. Each epoch has
125 channels at 500\,Hz (1{,}651 samples); as in the released code we keep the 122 EEG channels and
scale volts to microvolts, with no further normalisation. Every sentence maps to one of 39 semantic
categories through the released \texttt{textmaps.json}. Uniform chance is $1/39 = 2.56\%$; always
predicting the most frequent training category scores 4.72\% on the evaluation set described below.

\paragraph{The published protocol, as released.}
The paper describes the model (an EEGNet front end, one Transformer encoder layer, mean pooling,
a linear classifier) but not the split, the number of training runs, or how the reported accuracy is
chosen. We take these from the released code \citep{chiscocode} and the command its README gives
for ``reproducing the paper results'' (Appendix~\ref{app:upstream}). Trials from all 45 runs are
concatenated in recording order and every trial whose position satisfies $i \bmod 5 = 1$ is placed in
the evaluation set (1{,}249 trials); the other 4{,}994 are used for training. There is no third set.
Training runs for 100 epochs with AdamW (learning rate $5\times10^{-4}$, weight decay 0.01, batch 16,
linear warm-up over the first 10\% of steps, then linear decay). Every 100 optimiser steps the model
is evaluated on the evaluation set, and the script prints the maximum accuracy seen over all
evaluations: 313 of them in this configuration. We call this the \emph{max-over-checkpoints} rule.
Table~3 of the paper reports 13.88\% for S1 and a 5.52\% baseline obtained by shuffling the
correspondence between EEG and labels in training.

\paragraph{Our re-implementation.}
We copy the released model code and change only how its two temporal convolutions are computed
(an equivalent unfold-and-multiply form with the same parameters and initialisation, which is
about nine times faster on our hardware; outputs agree with the original to float32 rounding).
Data loading, label mapping and the split follow the released code and are checked by unit tests.
To read a second selection rule from the same runs, we hold out a random 10\% of the training
trials (499) as a validation set, so our models train on 4{,}495 trials and are evaluated
281 times. All runs use Apple-silicon GPUs (PyTorch MPS, float32) and three seeds.

\paragraph{Three readings of one run.}
For every run we store validation and evaluation-set accuracy at every evaluation, and the
evaluation-set predictions at three checkpoints:
\textbf{max-over-checkpoints} (the published rule);
\textbf{validation-selected}, the checkpoint with the highest validation accuracy, scored once on
the evaluation set, which is never used for selection; and
\textbf{last}, the final checkpoint.
Because all three come from the same training run, differences between them are due to the
selection rule alone.

\paragraph{Splits.}
Besides the released split (\emph{interleaved}) we use three others, each leaving about 80\% of
trials for training:
\emph{random-trial} (a random 20\% of trials);
\emph{run holdout} (the 45 runs in five fixed folds of nine runs; we report held-out folds 0--2); and
\emph{session holdout} (one recording day held out at a time; we report days 1--3).
For the grouped splits the validation set is made of whole training runs, so model selection is also
free of within-run neighbours. For cost reasons the split comparison uses the same model at
250\,Hz (pairwise averaging) and 30 epochs; it is compared only within itself. As a positive control
we also train on the reading epochs (the first 5\,s of each trial; 6{,}387 released epochs for S1,
2{,}501 samples at 500\,Hz), whose labels and trial order follow the same stimulus lists.

\paragraph{Permutation control.}
For each configuration we also train with the training labels randomly permuted; validation and
evaluation labels are untouched. This mirrors the shuffled baseline in the released code and gives
the accuracy each selection rule produces when the EEG carries no information about the label.

\paragraph{Uncertainty.}
The unit of replication for EEG models is a training seed ($n=3$). We report per-seed values, means
and standard deviations, and do not run significance tests on three seeds. For accuracies of a single
classifier on the evaluation set we report 95\% intervals from a bootstrap that resamples test
\emph{runs} (2{,}000 draws), because trials within a run share stimulus-list structure and are not
independent.
